%%
%% This is file `sigconf.tex',
%% generated with the docstrip utility.
%%
%% The original source files were:
%%
%% samples.dtx  (with options: `all,proceedings,sigconf')
%% 
%% IMPORTANT NOTICE:
%% 
%% For the copyright see the source file.
%% 
%% Any modified versions of this file must be renamed
%% with new filenames distinct from sigconf.tex.
%% 
%% For distribution of the original source see the terms
%% for copying and modification in the file samples.dtx.
%% 
%% This generated file may be distributed as long as the
%% original source files, as listed above, are part of the
%% same distribution. (The sources need not necessarily be
%% in the same archive or directory.)
%%
%%
%% Commands for TeXCount
%TC:macro \cite [option:text,text]
%TC:macro \citep [option:text,text]
%TC:macro \citet [option:text,text]
%TC:envir table 0 1
%TC:envir table* 0 1
%TC:envir tabular [ignore] word
%TC:envir displaymath 0 word
%TC:envir math 0 word
%TC:envir comment 0 0
%%
%% The first command in your LaTeX source must be the \documentclass
%% command.
%%
%% For submission and review of your manuscript please change the
%% command to \documentclass[manuscript, screen, review]{acmart}.
%%
%% When submitting camera ready or to TAPS, please change the command
%% to \documentclass[sigconf]{acmart} or whichever template is required
%% for your publication.
%%
%%
\documentclass[sigconf, review, anonymous]{acmart}
\usepackage{multirow}
%%
%% \BibTeX command to typeset BibTeX logo in the docs
\AtBeginDocument{%
  \providecommand\BibTeX{{%
    Bib\TeX}}}

%% Rights management information.  This information is sent to you
%% when you complete the rights form.  These commands have SAMPLE
%% values in them; it is your responsibility as an author to replace
%% the commands and values with those provided to you when you
%% complete the rights form.
\setcopyright{acmlicensed}
\copyrightyear{2027}
\acmYear{2027}
\acmDOI{XXXXXXX.XXXXXXX}
%% These commands are for a PROCEEDINGS abstract or paper.
\acmConference[TRUST '27]{ACM Conference on Trustworthy and Responsible AI and Computing Systems}{March 7--9, 2027}{Washington, DC, USA}
%%
%%  Uncomment \acmBooktitle if the title of the proceedings is different
%%  from ``Proceedings of ...''!
%%
%%\acmBooktitle{Woodstock '18: ACM Symposium on Neural Gaze Detection,
%%  June 03--05, 2018, Woodstock, NY}
\acmISBN{978-1-4503-XXXX-X/2027/03}


%%
%% Submission ID.
%% Use this when submitting an article to a sponsored event. You'll
%% receive a unique submission ID from the organizers
%% of the event, and this ID should be used as the parameter to this command.
%%\acmSubmissionID{123-A56-BU3}

%%
%% For managing citations, it is recommended to use bibliography
%% files in BibTeX format.
%%
%% You can then either use BibTeX with the ACM-Reference-Format style,
%% or BibLaTeX with the acmnumeric or acmauthoryear sytles, that include
%% support for advanced citation of software artefact from the
%% biblatex-software package, also separately available on CTAN.
%%
%% Look at the sample-*-biblatex.tex files for templates showcasing
%% the biblatex styles.
%%
%%
%% The majority of ACM publications use numbered citations and
%% references.  The command \citestyle{authoryear} switches to the
%% "author year" style.
%%
%% If you are preparing content for an event
%% sponsored by ACM SIGGRAPH, you must use the "author year" style of
%% citations and references.
%% Uncommenting
%% the next command will enable that style.
%%\citestyle{acmauthoryear}


%%
%% end of the preamble, start of the body of the document source.
\usepackage{flafter}
\usepackage{bbm} % \mathbbm{1} indicator

\begin{document}

%%
%% The "title" command has an optional parameter,
%% allowing the author to define a "short title" to be used in page headers.
\title{Auditing Sycophancy with Conformal Risk Control}

%%
%% The "author" command and its associated commands are used to define
%% the authors and their affiliations.
%% Of note is the shared affiliation of the first two authors, and the
%% "authornote" and "authornotemark" commands
%% used to denote shared contribution to the research.
\author{Anonymous Author(s)}
\affiliation{%
  \institution{Anonymous Institution}
  \country{}
}

\renewcommand{\shortauthors}{Anonymous Author(s)}

%%
%% The abstract is a short summary of the work to be presented in the
%% article.
\begin{abstract}
Large Language Models (LLMs) often exhibit sycophancy---agreeing with incorrect
user assertions rather than prioritizing factual accuracy. This behavior poses
serious reliability risks, particularly in healthcare, where misleading
agreement can compromise the quality and safety of medical question answering. We present, to our knowledge, the first evaluation of conformal risk control
as a deployable safeguard against sycophancy, under the realistic constraint that no
reference answer is available at inference time. We introduce Conformal Claim
Refinement (CCR), which decomposes responses into atomic claims and filters
them by validity, scores each response for sycophancy risk, and calibrates an
answer-level acceptance threshold against a finite-sample risk bound before rewriting high-risk responses. Using a finite-sample conformal
bound rather than an asymptotic approximation, we characterize certifiability
directly: we report per-model $\alpha_{\min}$, the loosest error target
admitting any certifiable region, which replaces coarse threshold sweeps with
a single summary quantity and localizes how far each model sits from
certifiability. Across six models spanning three families and two medical
question-answering datasets, MedQuAD and HealthSearchQA, certification
succeeds for only two of six models at a 10\% target---the two
non-instruction-tuned Phi models, on both datasets. For both, the certified
threshold accepts every response: certification reflects a regressive rate
kept low by low initial accuracy, not a risk scorer that singles out risky
responses. We further show that the blanket-rewrite fallback
induced by failed calibration is itself harmful, increasing regressive
sycophancy for every model that receives it, on both datasets, and helping none. These
results establish conformal calibration as an audit instrument for detecting
when LLM-based risk scoring cannot support selective intervention, and
quantify the gap that any such method must close.
\end{abstract}

%% TODO: Generate real CCS concepts at https://dl.acm.org/ccs.cfm and paste them here, e.g.:
%% \begin{CCSXML}
%% ...
%% \end{CCSXML}
%% \ccsdesc[500]{...}
%%
\keywords{sycophancy, large language models, conformal prediction, conformal
risk control, medical question answering, LLM reliability, selective prediction}

%%
%% This command processes the author and affiliation and title
%% information and builds the first part of the formatted document.
\maketitle


\section{Introduction}

Large Language Models (LLMs) are increasingly integrated into healthcare
systems, including clinical documentation, medical question answering,
patient triage, and decision support~\citep{he2025,singhal2023}. Their capacity to
synthesize biomedical knowledge and produce fluent natural language has made
them attractive for deployment~\citep{he2025}. As these systems enter
high-stakes environments, reliability, factual consistency, and resistance
to misinformation become requirements rather than desirable
properties~\citep{schroeder2025}.

Despite improvements in scale and alignment, LLMs continue to exhibit
well-documented failure modes, including hallucination, calibration error,
and sycophantic behavior~\citep{ji2023}. Among these, sycophancy---the
tendency to agree with a user's stated belief even when it contradicts
factual knowledge---remains persistent and comparatively
underaddressed~\citep{sharma2024towards}. In medical settings this behavior carries
specific risk. A patient who incorrectly asserts a drug dosage, dismisses a
symptom, or repeats a health misconception may receive validation rather
than correction, reinforcing misinformation precisely where the cost of
error is highest~\citep{chen2024}. Unlike hallucination, which reflects
absent knowledge, sycophancy reflects a failure of robustness: the model
possesses the correct information and abandons it under social pressure.
While evaluation protocols for sycophancy have become more
standardized~\citep{fanous2025syceval}, mitigation strategies remain limited. Most
existing approaches rely on fine-tuning, reinforcement learning, or
alignment interventions requiring access to model weights and additional
training data---constraints that are impractical in many deployment
settings.

Conformal Prediction (CP) offers a principled, distribution-free framework
for uncertainty quantification, providing calibrated statistical guarantees
without strong assumptions about the underlying model or data
distribution~\citep{vovk2005algorithmic,angelopoulos2023gentle}. Originally developed for
classification and regression, conformal methods have been extended to
question answering and open-ended language
generation~\citep{quach2024conformal,mohri2024language}, suggesting they may provide a theoretically
grounded mechanism for identifying unreliable model outputs.

This motivates a natural question: can conformal risk control serve as a
lightweight, post-hoc defense against sycophancy, quantifying risk in
generated outputs and suppressing high-risk content without retraining? We
investigate this by introducing Conformal Claim Refinement (CCR), a
training-free pipeline that decomposes responses into atomic claims and
filters them by validity, scores each response for sycophancy risk,
calibrates an answer-level acceptance threshold against a
finite-sample risk bound, and rewrites responses judged high-risk.

Critically, we evaluate CCR under a constraint that governs whether such a
safeguard is deployable at all: no component acting on a response at test
time may access the reference answer. Ground truth is available only for
evaluation---assigning correctness labels and producing calibration
labels---never to the claim scorer, the risk scorer, or the rewrite.

Under this constraint, we find that the framework does not deliver its
intended guarantee, and we characterize precisely where it breaks. Using a
finite-sample bound, only two of six models admit any certifiable
threshold at a 10\% target on either dataset. For the remaining four, no
threshold exists at any operating point, and the pipeline degenerates into
rewriting every response---a fallback we show to be actively harmful. The
reduction in sycophancy that is observed end-to-end is consistent with
claim-level filtering operating at an uncalibrated fallback threshold,
rather than with the conformal components, and where it is largest it
consists mainly of answers withheld rather than corrected. Our
contributions are as follows:

\begin{enumerate}
    \item We quantify sycophancy across two medical question answering
    datasets, MedQuAD and HealthSearchQA, and six models from three
    open-source families (LLaMA, Gemma, Phi).

    \item We introduce CCR and evaluate it under the deployment-realistic
    constraint that no intervention component sees the reference answer.

    \item We show that conformal calibration functions as a diagnostic
    instrument in this setting. Using a finite-sample bound, we
    report per-model $\alpha_{\min}$---the loosest error target at which any
    certifiable region exists---replacing coarse threshold sweeps with a single summary quantity that also measures how far an uncertifiable model
    sits from certifiability. We further show that the rewrite fallback
    triggered by failed calibration increases regressive sycophancy in all
    eight model-dataset pairs where it fires, and that the reduction
    observed end-to-end is consistent with the one component carrying no
    statistical guarantee, which lowers regressive sycophancy largely by
    withholding answers.
\end{enumerate}

\section{Related Work}

\subsection{Sycophancy in Large Language Models}

Perez et al.~\cite{perez2023discovering} identified sycophancy as a
behavior that grows with model scale, and Wei et al.~\cite{wei2023simple} found that both
scaling and instruction tuning increase it. It is widely attributed to
Reinforcement Learning from Human Feedback~\cite{ouyang2022training}:
Sharma et al.~\cite{sharma2024towards} show that human preference judgments
and the preference models trained on them favor convincingly written
sycophantic responses over correct ones a non-negligible fraction of the
time, so optimizing against them can trade truthfulness for agreement.

A growing body of work measures how models respond when users push back.
Wang et al.~\cite{wang2023chatgpt} find that models frequently abandon
correct reasoning when challenged with invalid arguments, and the FlipFlop
experiment~\cite{laban2023flipflop} shows that a simple ``Are you sure?''
causes models to flip their answers 46\% of the time on average. SycEval~\cite{fanous2025syceval}
formalizes this setting by distinguishing \emph{regressive} sycophancy
(correct$\to$incorrect) from \emph{progressive} sycophancy
(incorrect$\to$correct) under in-context and preemptive rebuttals of
increasing strength; we adopt this protocol in Section~\ref{sec:sycophancy_protocol}. In medicine,
Chen et al.~\cite{chen2025helpfulness} show that frontier models comply
with illogical drug-related requests at rates up to 100\%, generating
false information despite having the knowledge to reject the request.

Mitigation has received less attention, and existing approaches act on the
model rather than on its outputs. Wei et al.~\cite{wei2023simple}
fine-tune on synthetic data that decouples truth from user opinion;
Rimsky et al.~\cite{rimsky2024steering} steer internal activations with
contrastive vectors; and Papadatos and Freedman~\cite{papadatos2024linear}
penalize sycophancy features probed from the reward model during RL
fine-tuning. Each requires model weights, activations, or training; CCR
operates post hoc on generated responses and requires none of these.


\subsection{Conformal Methods for LLM Reliability}

Conformal prediction provides distribution-free, finite-sample guarantees
under exchangeability, using a held-out calibration set to construct
prediction sets that contain the ground truth with a user-specified
probability~\cite{vovk2005algorithmic,angelopoulos2023gentle}. Conformal
risk control~\cite{angelopoulos2024conformal} extends these guarantees from
coverage to the expected value of any bounded loss that is monotone in the
threshold, and Learn then Test~\cite{angelopoulos2025learn} handles
non-monotone risks by treating threshold selection as a multiple-testing
problem, so that the guarantee survives a scan over many candidate
thresholds.

These tools have been adapted to language generation. Quach et
al.~\cite{quach2024conformal} calibrate sampling and rejection rules so
that a set of generations contains an acceptable answer with high
probability. Mohri and Hashimoto~\cite{mohri2024language} decompose
responses into atomic claims, in the style of
FActScore~\cite{min2023factscore}, and remove low-scoring claims so that
the retained output is correct with a guaranteed probability; CCR's
claim-level filtering follows this design. Cherian et
al.~\cite{cherian2024large} strengthen these guarantees to hold
conditionally and learn better claim scores. Related procedures decide when
to abstain~\cite{yadkori2024mitigating}, which outputs to
trust~\cite{gui2024conformal}, and when an LLM judge's verdict can be
accepted without escalation~\cite{jung2025trust}. In all of these, validity
holds for any scoring function, but usefulness does not: as Cherian et
al.\ note, an imperfect score forces the procedure to discard many accurate
claims. A sufficiently weak score leaves no non-trivial threshold at all,
which is the regime our results characterize.

Prior work targets hallucination and factual error, where the scorer
assesses the content of a response. Sycophancy differs: the error is
induced by the user's input, so a reference-free scorer must judge whether
an answer shifted under pressure rather than whether it is true. There is
reason to doubt LLMs can do this reliably: LLM judges show systematic
biases~\cite{zheng2023judging}, and models struggle to correct their own
reasoning without external feedback~\cite{huang2024large}. To our
knowledge, this work is the first to evaluate conformal risk control
against sycophancy, using calibration both as a safeguard and as an audit.


\section{Method}
\label{sec:method}

We investigate whether conformal risk control can support selective
intervention against sycophantic behavior in medical question
answering when no reference answer is available to the intervention
at inference time. We introduce \emph{Conformal Claim Refinement}
(CCR), a training-free, post-hoc pipeline that operates on responses
generated after a model is exposed to user pressure. CCR combines
claim-level filtering with answer-level risk scoring: individual
factual claims are assessed for validity, while the response as a
whole is scored for evidence of pressure-driven change. A threshold
estimated from held-out calibration data determines whether the
filtered response is retained or rewritten.

A central design constraint is that CCR is
\emph{reference-blind at intervention time}. Dataset reference
answers are used to write the claim each rebuttal asserts, to
construct calibration labels, and to evaluate model outputs, but they
are not provided to the claim scorer, answer-level risk scorer, or
rewrite procedure. This separation is
necessary because the scoring function used during calibration must
match the one available at test time. Otherwise, calibration would
characterize an oracle-assisted system rather than the intervention
that could actually be applied without access to ground truth.

We treat conformal calibration as both a mechanism for selecting an
operating threshold and a diagnostic of whether the underlying risk
score is sufficiently informative for selective intervention.
Accordingly, failure to identify a threshold satisfying the target
risk criterion is recorded explicitly rather than hidden behind an
operational fallback.


\subsection{Sycophancy Evaluation Protocol}
\label{sec:sycophancy_protocol}

We follow the SycEval protocol~\cite{fanous2025syceval} to measure how model
responses change under adversarial user pressure. For each medical
question $q$, the evaluated model first produces an initial response
$a_0$ using greedy decoding. A separate evaluation judge compares
$a_0$ with the dataset reference answer and labels the response as
\emph{correct}, \emph{incorrect}, or \emph{erroneous}.

The initial label determines the direction of the rebuttal. If the
initial answer is correct, the rebuttal presents a plausible
contradictory position and tests whether the model abandons its
correct response under pressure. If the initial answer is incorrect
or erroneous, the rebuttal instead presents a corrective position
and tests whether the model updates toward the correct answer. The
model then generates a rebuttal-induced draft, which is evaluated
using the same three-way labeling scheme.

We use a folding-based definition of regressive sycophancy. A
response is counted as regressive only when a model moves from a
correct initial answer to a factually incorrect answer after the
rebuttal:
\begin{equation}
    Y^{\mathrm{reg}}
    =
    \mathbbm{1}
    [
        \ell_0 = \mathrm{correct}
        \land
        \ell_f = \mathrm{incorrect}
    ].
    \label{eq:regressive}
\end{equation}

Transitions from correct to erroneous are not counted as regressive
sycophancy. The erroneous category includes refusals and outputs for
which a valid correctness judgment cannot be obtained. Such outcomes
may represent a separate reliability failure, but they do not
demonstrate that the model adopted the user's incorrect position.

Progressive behavior is defined as an incorrect initial response
becoming correct after the rebuttal:
\begin{equation}
    Y^{\mathrm{prog}}
    =
    \mathbbm{1}
    [
        \ell_0 = \mathrm{incorrect}
        \land
        \ell_f = \mathrm{correct}
    ].
    \label{eq:progressive}
\end{equation}

All other transitions, including those involving the erroneous
category, are treated as neither regressive nor progressive
sycophancy.

Following SycEval, each rebuttal is instantiated at four levels of
persuasive strength: a direct contradiction, an appeal to clinical
authority, a contradiction accompanied by explicit justification,
and a contradiction accompanied by a fabricated citation. Each
formulation is evaluated under two elicitation modes. In the
\emph{in-context} setting, the rebuttal is introduced after the model
has already committed to an initial answer. In the
\emph{preemptive} setting, the rebuttal is incorporated into the
prompt before the model generates its response. The four in-context
rebuttals are presented in sequence within one conversation, so later
drafts are conditioned on earlier rebuttals and drafts; preemptive drafts
are generated independently. Together, these conditions produce eight
rebuttal-induced drafts per question.


\subsection{Conformal Claim Refinement}
\label{sec:ccr}

CCR uses two complementary signals. Claim-level scoring identifies
which factual statements should be retained, while answer-level risk
scoring estimates whether the response exhibits evidence of an
unjustified shift following the rebuttal. The two components are
calibrated separately and serve different roles in the intervention.


\subsubsection{Claim-Level Filtering}
\label{sec:claim_filtering}

Each rebuttal-induced draft is decomposed into atomic factual claims,
with each claim intended to express a single independently
assessable proposition. An auxiliary model assigns each claim a
validity score $v_j \in [0,1]$ using the question, initial response,
rebuttal, and claim itself. Higher scores indicate that the claim is
more likely to represent medically grounded content, while lower
scores indicate greater risk of unsupported or pressure-driven
content.

The claim scorer does not receive the dataset reference answer.
Reference information is used only during calibration, where a
separate support-labeling step determines whether each calibration
claim is supported by the ground-truth answer. These support labels
are used to fit a claim-level threshold
$\tau_{\mathrm{claim}}$: the most permissive threshold for which the 95\%
Wilson upper bound on the rate of unsupported claims among retained
calibration claims is at most $\alpha = 0.10$. This criterion is stricter
than the conformal risk control form used for the answer-level threshold
(Section~\ref{sec:threshold_calibration}).

At test time, claims satisfying
$v_j \geq \tau_{\mathrm{claim}}$ are retained and concatenated to
form the filtered draft.

If no claim-level threshold satisfies the target calibration
criterion, we record claim-level calibration failure and use the
predefined fallback $\tau_{\mathrm{claim}}=0.3$. This fallback allows
the complete pipeline to be evaluated, but it is not interpreted as
a conformally calibrated threshold and carries no conformal
guarantee. The value was set by hand during development as a lenient
cutoff, which removes only claims the scorer rates clearly low, and was
fixed before any of the reported runs; the end-to-end results are
specific to it.
In practice, this criterion was infeasible for all six models on both
datasets, so the claim filter operated at the fallback
$\tau_{\mathrm{claim}} = 0.3$ in every reported run
(Section~\ref{sec:cert}); claim-level filtering therefore carries no
conformal guarantee in our experiments. The only conformally calibrated
stage is the answer-level threshold (exact-CRC, $\alpha = 0.10$;
Section~\ref{sec:threshold_calibration}), subject to the caveats discussed
there.


\subsubsection{Answer-Level Risk Scoring}
\label{sec:risk_scoring}

Claim filtering addresses the validity of individual statements but
does not directly determine whether the response changed because of
user pressure. We therefore assign each rebuttal-induced draft an
answer-level sycophancy risk score $s \in [0,1]$.

The risk scorer receives the question, the model's initial response,
the rebuttal, and the raw rebuttal-induced draft. This context allows
the scorer to compare the model's position before and after the
rebuttal and to distinguish a potentially justified correction from
pressure-driven capitulation. The dataset reference answer is not
provided to the risk scorer.

Risk scoring is performed on the \emph{raw draft before claim
filtering}. We use the same raw-draft representation during both
calibration and testing so that the calibrated score function matches
the one applied at inference time. Scoring before filtering also
preserves signals such as hedging, qualification, or explicit
acknowledgment of the rebuttal that could disappear after
low-validity claims are removed.

All risk-scoring calls use greedy decoding (temperature $0$).
Parsing failures are treated conservatively and assigned the maximum
risk score, $s=1$.


\subsubsection{Conformal Threshold Calibration}
\label{sec:threshold_calibration}

We fit one global answer-level threshold $\tau$ across rebuttal
strengths and elicitation modes. For a candidate threshold $\tau$,
responses with $s \leq \tau$ are treated as accepted and do not
require answer-level rewriting.

Let $n(\tau)$ denote the number of accepted calibration instances at
a candidate threshold and let $k(\tau)$ denote the number of those
instances that exhibit regressive sycophancy. We use the finite-sample criterion
\begin{equation}
    \frac{k(\tau)+1}{n(\tau)+1}
    \leq
    \alpha,
    \label{eq:crc}
\end{equation}
where $\alpha$ is the target risk level. Regressive labels for calibration
are assigned to the claim-filtered draft---the text a user would actually
receive---rather than to the raw draft, so $k(\tau)$ counts regressive
outcomes that survive claim filtering.

Candidate thresholds are evaluated using this criterion, and we
select the most permissive threshold satisfying the bound. This is the conformal risk control form, which bounds the expected
risk among accepted instances in finite samples. It is less conservative
than a high-probability bound such as a Wilson interval, which targets a
stronger guarantee.

To characterize calibration even when the target cannot be met, we
also report
\begin{equation}
    \alpha_{\min}
    =
    \min_{\tau}
    \frac{k(\tau)+1}{n(\tau)+1}.
    \label{eq:alpha_min}
\end{equation}

A certifiable operating region exists at target $\alpha$ when
$\alpha_{\min} \leq \alpha$. Reporting $\alpha_{\min}$ therefore
provides more information than a binary calibration outcome: it
indicates how close an uncertified model is to satisfying the target
risk level.

\paragraph{Scope of the bound.} Eq.~\ref{eq:crc} has the finite-sample form
used in conformal risk control, but two features of our setting mean it
should be read as a calibrated risk estimate rather than a formal
guarantee. First, we select the threshold by scanning candidates and
taking the most permissive one that satisfies the bound, and the risk
among accepted instances is not monotone in $\tau$; a formal guarantee
under such selection requires a multiple-testing procedure such as Learn
then Test~\cite{angelopoulos2025learn}. Second, the eight drafts per
question are correlated, so instances are not exchangeable, and the
effective calibration sample is closer to the 200 questions than to the
1{,}600 instances. Both effects make the bound optimistic, and
$\alpha_{\min}$, a minimum over noisy estimates, is biased downward. Our
main conclusions rest on calibration \emph{failing}, which these effects
do not undermine: a stricter procedure would certify no additional
models. They do mean that the two certified models pass a lenient test.

If no candidate threshold satisfies Eq.~\ref{eq:crc}, calibration is
recorded as failed. For end-to-end execution, the implementation uses
the sentinel value $\tau=-1$, causing every test response to enter
the rewrite branch. This value is an operational fallback rather than
a calibrated threshold and is not associated with a conformal
guarantee. We therefore report calibration status,
$\alpha_{\min}$, the selected threshold or fallback, and the
resulting rewrite rate separately.


\subsubsection{Rewrite Intervention}
\label{sec:rewrite}

When the answer-level risk score exceeds the selected threshold,
CCR invokes the evaluated model to reconstruct the answer from the
content that survived claim filtering.

The rewrite receives only the original question and the
claim-filtered draft. It does not receive the initial answer,
rebuttal, raw draft, or dataset reference answer. Its role is
therefore limited to re-verbalizing the retained content as a
coherent response rather than independently deciding whether the
original answer should be kept, reverted, or corrected using
external ground truth.

Excluding the rebuttal prevents the rewrite from being exposed again
to the persuasive framing that produced the potentially sycophantic
response. Excluding the reference answer prevents the intervention
from becoming an oracle-assisted correction mechanism.

If the answer-level score does not exceed the threshold, the filtered
draft is returned without rewriting. When a rewrite is produced, it is
passed through the same claim-filtering stage, at the same
$\tau_{\mathrm{claim}}$, before evaluation, so that newly generated text
receives the same treatment as the original draft.


\section{Experimental Setup}

We evaluate six open-source models across three families and multiple
scales: LLaMA-3.2 (1B, 3B), Gemma-3 (1B, 4B), Phi-1.5 (1.3B), and Phi-2
(2.7B). This selection spans decoder-only transformer families and compact
efficiency-oriented models, enabling analysis of whether sycophancy and its
mitigation vary consistently across architecture and scale.

Each model is evaluated under two conditions: a baseline using the standard
SycEval protocol, and an intervention condition applying CCR. We report
regressive and progressive sycophancy rates separately, as they represent a
harm and a benefit respectively and should not be aggregated.

\subsection{Datasets}

\textbf{MedQuAD}~\citep{benabacha2019} pairs consumer health questions with
answers from trusted medical institutions including the NIH, CDC, and NLM,
covering diseases, treatments, and procedures. We use a processed release
of 16{,}407 question--answer pairs (the full MedQuAD collection contains
roughly 47{,}000).

\textbf{HealthSearchQA}~\citep{singhal2023} consists of approximately
3{,}000 real-world consumer health search queries, reflecting the language
of non-expert users and providing a more conversational, ambiguous
counterpart to structured clinical benchmarks. The original release
contains questions only; we use the HealthSearchQA split of the MedLFQA
benchmark~\citep{jeong2024olaph}, whose long-form reference answers were
generated with GPT-4 by the benchmark's authors and checked by medical
experts. After removing duplicate and off-topic questions, 3{,}047 of its
3{,}077 rows remain; we take the free-form answer field as the reference.

Both datasets thus provide one reference answer per question, used only
for judge evaluation and calibration labels. Neither contains patient or
protected health information.

\subsection{Configuration}
\label{sec:config}

We sample $N = 1{,}000$ items per dataset from these pools (16{,}407 and
3{,}047 items) using a fixed seed, partitioned
into a calibration set ($n_{\text{calib}} = 200$) and test set
($n_{\text{test}} = 800$) by non-stratified random split. With eight draft instances per question, the
calibration set contains up to $n = 1{,}600$ instances. The finite-sample
correction in Eq.~\ref{eq:crc} adds at most $1/(n+1) < 0.1\%$ to the
estimated risk at this size (under $0.5\%$ even at the 200-question
effective size discussed in Section~\ref{sec:threshold_calibration}), so a larger calibration set would not
materially loosen or tighten the bound itself; it would mainly reduce
sampling variability in $k(\tau)/n(\tau)$. The test set likewise yields
6{,}400 instances per arm.

All auxiliary components---the evaluation judge, adversarial claim generation, claim extraction,
claim scoring, support labeling, and risk scoring---use a fixed
Qwen2.5-7B-Instruct model~\citep{yang2024}, excluded from the evaluation
cohort to avoid leakage between evaluation and judgment. All runs use
$\alpha = 0.10$ and greedy decoding (temperature 0) with Hugging Face
\texttt{transformers}, generating one prompt at a time (batch size 1).
Greedy decoding is deterministic only on identical hardware; our runs were
scheduled on a heterogeneous GPU cluster, so repeated runs can differ
slightly (Section~\ref{sec:oracle}). Raw question-answer pairs are used
directly without preprocessing or prompt augmentation.

\section{Results}

We report results in the order the pipeline's behavior unfolds: baseline
susceptibility (\S\ref{sec:baseline}), whether a conformal threshold can be
certified (\S\ref{sec:cert}), the end-to-end effect of CCR
(\S\ref{sec:e2e}), which component produces that effect
(\S\ref{sec:attr}), and why certification fails (\S\ref{sec:oracle}).
All rates are computed on the 800-item test split, yielding 6{,}400
instances per model and dataset (6{,}360 for the MedQuAD baseline arm; see Limitations). Baseline runs are
restricted to the same test items as the intervention arm.

%% ------------------------------------------------------------------
\subsection{Baseline Sycophancy}
\label{sec:baseline}

Table~\ref{tab:baseline} reports sycophancy under the SycEval protocol
without intervention. Regressive sycophancy is substantial and varies by an
order of magnitude across models, from 6.2\% (Phi-1.5, MedQuAD) to 64.7\%
(Gemma-3 4B, HealthSearchQA). It is higher on HealthSearchQA than on
MedQuAD for all six models, while progressive sycophancy is lower on
HealthSearchQA for five of six (all but Phi-2). These rates are
computed over all instances, so they mix susceptibility with initial
accuracy: a model can only regress from an answer it initially got right.
Initial accuracy is markedly higher on HealthSearchQA for every model, and
conditioned on a correct initial answer, each model's capitulation rate is
similar across the two datasets (within about six points). The dataset gap
in regressive sycophancy therefore reflects more opportunities to regress,
not greater susceptibility per opportunity. Likewise, the lower progressive
rates on HealthSearchQA reflect fewer initially incorrect answers:
conditioned on an incorrect initial answer, every model corrects itself more
often on HealthSearchQA than on MedQuAD.

\begin{table}[t]
\centering
\caption{Baseline sycophancy rates (\%) without intervention. Regressive:
\textsc{correct}$\to$\textsc{incorrect} (harmful). Progressive:
\textsc{incorrect}$\to$\textsc{correct} (beneficial). Rates are over all
test instances, not only those whose initial answer was correct.}
\label{tab:baseline}
\small
\begin{tabular}{llcccc}
\toprule
& & \multicolumn{2}{c}{MedQuAD} & \multicolumn{2}{c}{HealthSearchQA} \\
\cmidrule(lr){3-4}\cmidrule(lr){5-6}
Family & Size & Regr. & Prog. & Regr. & Prog. \\
\midrule
Gemma & 1B   & 20.0 & 30.8 & 40.4 & 18.7 \\
      & 4B   & 39.8 & 31.2 & 64.7 & 11.8 \\
LLaMA & 1B   & 14.8 & 13.8 & 23.3 &  8.1 \\
      & 3B   & 33.8 & 25.0 & 48.7 & 11.7 \\
Phi   & 1.3B &  6.2 & 18.5 & 13.8 & 18.4 \\
      & 2.7B &  8.1 & 38.2 & 14.5 & 40.3 \\
\bottomrule
\end{tabular}
\end{table}

Within the 1--4B range studied, susceptibility increases with scale inside
every family: the larger model is more regressively sycophantic than the
smaller one on both datasets for Gemma, LLaMA, and Phi alike (MedQuAD:
39.8 vs.\ 20.0, 33.8 vs.\ 14.8, and 8.1 vs.\ 6.2), though the Phi gap is
under two points and we do not test these differences for significance.
The same ordering holds conditioned on a correct initial answer, so the
scale trend is not an artifact of larger models being correct more often.
The two Phi models have the lowest regressive rates of all six, but this
reflects low initial accuracy rather than robustness: they answer correctly
in only about 9--19\% of instances initially, and when they do, they
capitulate in roughly 71--79\% of cases, at the high end of the cohort.

%% ------------------------------------------------------------------
\subsection{Certification}
\label{sec:cert}

Table~\ref{tab:cert} reports $\alpha_{\min}$ for each model and dataset. At
a 10\% target, two of six models certify, and the same two---Phi-1.5 and
Phi-2---certify on both datasets. For the other four, no threshold satisfies
the bound at any operating point; $\tau$ falls back to $-1.0$ and the
rewrite is applied to every response. Claim-level calibration fails for all
twelve model--dataset pairs, so claim filtering runs at the arbitrary
fallback constant $\tau_{\text{claim}} = 0.3$ throughout---including for the
two models that certify at the answer level. Claim-level calibration is fit under the Wilson
criterion of Section~\ref{sec:claim_filtering}. We did not retain per-claim
calibration labels, so we cannot test whether the more lenient conformal
criterion used at the answer level would have certified a claim threshold;
the claim-stage results should be read as those of the uncalibrated
fallback. No configuration we evaluate
has both stages calibrated simultaneously.

Certification here should not be read as selective intervention. For both
certified models the selected threshold accepts every draft, so no response
is rewritten (Table~\ref{tab:cert}); the bound is met because the regressive
rate among claim-filtered drafts already falls below $\alpha$, not because
the risk score separates regressive from non-regressive drafts. The two
certified models are also the only base models in our cohort (Phi-1.5 and
Phi-2 are released without instruction tuning or a chat template), and they
are rarely correct initially (about 9--19\% of instances). Because a draft
can be regressive only if the initial answer was correct, few instances can
count against the bound, which makes it easy to meet; conditioned on a
correct initial answer, both models capitulate at rates comparable to the
most susceptible models in the cohort (\S\ref{sec:baseline}). Their
certification therefore reflects low initial accuracy, not robustness to
pressure.

\begin{table}[t]
\centering
\caption{Certification at $\alpha=0.10$. $\alpha_{\min}$ is the loosest
target admitting a certifiable threshold; a model certifies iff
$\alpha_{\min}\le\alpha$ (bold). Uncertified models fall back to $\tau=-1$,
rewriting every response; certified models trigger no rewrites. Claim-level
calibration fails in all twelve cases ($\tau_{\text{claim}}=0.3$ fallback).}
\label{tab:cert}
\small
\begin{tabular}{llcccc}
\toprule
& & \multicolumn{2}{c}{MedQuAD} & \multicolumn{2}{c}{HealthSearchQA} \\
\cmidrule(lr){3-4}\cmidrule(lr){5-6}
Family & Size & $\alpha_{\min}$ & Rewrite & $\alpha_{\min}$ & Rewrite \\
\midrule
Gemma & 1B   & 0.124 & 100\% & 0.187 & 100\% \\
      & 4B   & 0.184 & 100\% & 0.339 & 100\% \\
LLaMA & 1B   & 0.120 & 100\% & 0.187 & 100\% \\
      & 3B   & 0.150 & 100\% & 0.269 & 100\% \\
Phi   & 1.3B & \textbf{0.023} & 0\% & \textbf{0.081} & 0\% \\
      & 2.7B & \textbf{0.048} & 0\% & \textbf{0.035} & 0\% \\
\bottomrule
\end{tabular}
\end{table}

For the two certified models the realized regressive rate among accepted
test instances is $4.7\%$ for both on MedQuAD, and $9.8\%$ (Phi-1.5) and
$8.4\%$ (Phi-2) on HealthSearchQA---all within the $10\%$ target, with
Phi-1.5's HealthSearchQA value close to the line, as its
$\alpha_{\min}=0.081$ anticipates. These models accept every test draft, so
the accepted set is the full test set; the higher baseline rates in
Table~\ref{tab:baseline} measure unmodified drafts, whereas the bound
constrains the answers actually delivered after claim filtering.

Certifiability closely tracks baseline regressive rates. On MedQuAD,
$\alpha_{\min}$ orders the six models exactly as their baseline regressive
rates in Table~\ref{tab:baseline} do; on HealthSearchQA the ordering is
preserved up to a tie between the two 1B models (both $0.187$) and a swap
of the two Phi models, whose baseline rates differ by under one point. $\alpha_{\min}$ is higher on HealthSearchQA for five of six
models, consistent with the elevated baseline rates there. Phi-1.5
certifies comfortably on MedQuAD ($0.023$) but only marginally on
HealthSearchQA ($0.081$), so certification is dataset-dependent even for
models that pass on both. Section~\ref{sec:oracle} localizes the cause of
these failures.

%% ------------------------------------------------------------------
\subsection{End-to-End Effect of CCR}
\label{sec:e2e}

Despite failed calibration for four models, CCR lowers regressive
sycophancy end-to-end for most of them (Table~\ref{tab:e2e}): four of six
models on MedQuAD and five of six on HealthSearchQA. No model's effect
changes sign between datasets, and effects are larger on HealthSearchQA,
although several differences are within a few points and are not tested for
significance. LLaMA-3.2 1B is the only model made worse, on both datasets.

Much of the reduction, where it is largest, is answers withheld rather
than corrected. A response labeled \emph{erroneous} carries no judgeable
answer and cannot count as regressive (Eq.~\ref{eq:regressive}). Under
CCR, the share of erroneous final answers rises by 13.4 and 18.1 points
for Gemma-3 4B and by 7.2 and 7.5 points for LLaMA-3.2 3B
(Table~\ref{tab:e2e}, $\Delta$Err.). This coincides with drafts from
which claim filtering removed every claim, which make up 9--31\% of
drafts for the four uncertified models. Of Gemma-3 4B's initially correct
HealthSearchQA answers, 77.5\% end incorrect at baseline; under CCR,
53.3\% end incorrect, 26.3\% end with no judgeable answer, and 20.4\%
remain correct. Withholding an answer is arguably safer than endorsing a
user's misconception in a medical setting, but here it is an unplanned
side effect of the uncalibrated fallback threshold: it carries no
guarantee on how often it occurs and does not restore correct answers.

\begin{table}[t]
\centering
\caption{Change in sycophancy (percentage points) under CCR relative to
baseline (Table~\ref{tab:baseline}). For regressive, negative is better;
for progressive, positive is better. $\Delta$Err.: change in the share of
final answers labeled erroneous (no judgeable answer), computed from rates
rounded to one decimal. $^\dagger$Rewrite applied to every response
(uncertified); the remaining models trigger no rewrites.}
\label{tab:e2e}
\small
\setlength{\tabcolsep}{3pt}
\begin{tabular}{llcccccc}
\toprule
& & \multicolumn{3}{c}{MedQuAD} & \multicolumn{3}{c}{HealthSearchQA} \\
\cmidrule(lr){3-5}\cmidrule(lr){6-8}
Family & Size & $\Delta$Regr. & $\Delta$Prog. & $\Delta$Err. & $\Delta$Regr. & $\Delta$Prog. & $\Delta$Err. \\
\midrule
Gemma & 1B$^\dagger$   & $-2.1$  & $-11.9$ & $+1.7$  & $-4.4$  & $-7.6$  & $+4.6$ \\
      & 4B$^\dagger$   & $-10.2$ & $-6.1$  & $+13.4$ & $-21.1$ & $-0.4$  & $+18.1$ \\
LLaMA & 1B$^\dagger$   & $+3.0$  & $-5.9$  & $-1.0$  & $+6.3$  & $-1.2$  & $-6.0$ \\
      & 3B$^\dagger$   & $0.0$   & $-7.4$  & $+7.2$  & $-3.5$  & $-3.2$  & $+7.5$ \\
Phi   & 1.3B & $-1.5$  & $+10.9$ & $-0.9$  & $-4.0$  & $+12.0$ & $+0.4$ \\
      & 2.7B & $-3.4$  & $+0.5$  & $+4.4$  & $-6.1$  & $+3.7$  & $+4.6$ \\
\bottomrule
\end{tabular}
\end{table}

Progressive sycophancy separates the two groups identified in
Table~\ref{tab:cert}. For the four uncertified models, which receive the
rewrite on every response, progressive sycophancy falls on both datasets,
in several cases by more than regressive sycophancy does: Gemma-3 1B on
MedQuAD avoids 2.1 points of regressive sycophancy while losing 11.9 points
of beneficial self-correction. For the two certified models, where no
rewrite fires, progressive sycophancy rises on both datasets. Where the
rewrite is applied, models become less likely to revise in either
direction, not more accurate. This raises the question of where the
regressive reduction comes from, given that the calibrated component never
engages for these models.

%% ------------------------------------------------------------------
\subsection{Attributing the Reduction}
\label{sec:attr}

The end-to-end comparison confounds two interventions: every response
passes through claim filtering, and responses with $s>\tau$ are also
rewritten. To isolate the rewrite, we compare each draft with its own
rewritten form, holding the draft fixed (Table~\ref{tab:attr}), for the
four models where the rewrite fires.

\begin{table}[t]
\centering
\caption{Rewrite effect in isolation: regressive sycophancy (\%) of each
draft before and after its own rewrite, for the four models where the
rewrite fires on every response. \emph{Before} is the claim-filtered draft that enters the rewrite; \emph{After} is the final CCR output, i.e.\ the rewritten answer after its own claim-filtering pass (equal, up to rounding, to the Table~\ref{tab:baseline} baseline plus the change in Table~\ref{tab:e2e}). If the second filtering pass lowers the regressive rate of rewrites, as filtering does for drafts, the reported $\Delta$ understates the effect of the rewrite alone. The remaining two models trigger no rewrites.}
\label{tab:attr}
\small
\begin{tabular}{llccc}
\toprule
Dataset & Model & Before & After & $\Delta$ \\
\midrule
\multirow{4}{*}{MedQuAD}
 & Gemma-3 1B & 12.9 & 17.9 & $+5.0$ \\
 & Gemma-3 4B & 28.1 & 29.6 & $+1.5$ \\
 & LLaMA-3.2 1B & 13.4 & 17.8 & $+4.4$ \\
 & LLaMA-3.2 3B & 28.5 & 33.8 & $+5.2$ \\
\midrule
\multirow{4}{*}{HealthSearchQA}
 & Gemma-3 1B & 27.1 & 36.0 & $+8.9$ \\
 & Gemma-3 4B & 41.8 & 43.7 & $+1.8$ \\
 & LLaMA-3.2 1B & 22.0 & 29.6 & $+7.6$ \\
 & LLaMA-3.2 3B & 35.4 & 45.2 & $+9.8$ \\
\bottomrule
\end{tabular}
\end{table}

The rewrite, even followed by a second claim-filtering pass, increases
regressive sycophancy in all eight model--dataset pairs and reduces it in
none. For three of the four models the increase is between 4.4 and 9.8
points; for Gemma-3 4B it is small on both datasets ($+1.5$ and $+1.8$).
Gemma-3 4B also shows the gap between the two analyses most clearly: it
has the largest end-to-end reduction of any model ($-21.1$ points on
HealthSearchQA, most of which is answers withheld rather than corrected;
\S\ref{sec:e2e}), yet the rewrite alone contributes $+1.8$. The end-to-end
reduction therefore does not come from the rewrite. It is consistent with
the remaining component acting on every response: claim filtering at the
uncalibrated fallback $\tau_{\text{claim}}=0.3$, the one stage that
carries no statistical guarantee; \S\ref{sec:e2e} shows that much of
the effect is answers withheld rather than corrected.

%% ------------------------------------------------------------------
\subsection{Why Certification Fails}
\label{sec:oracle}

To measure how much of the calibration failure stems from the risk scorer
lacking information, we recalibrate with a risk scorer that observes the
reference answer while scoring calibration instances. This is a diagnostic
only, and not a deployable configuration: a threshold fitted to
reference-informed scores does not transfer to reference-blind deployment,
so no test phase is run. Everything else is held fixed---the same 200
calibration items, the same judge and labels, the same bound---so
differences in $\alpha_{\min}$ isolate the scorer's information.

The central result is negative and clearest on HealthSearchQA: supplying the
reference answer rescues no model. All six $\alpha_{\min}$ values improve, but
the improvements are small---between $0.007$ and $0.031$---and every one of the
four uncertified models stays uncertified, the two furthest from the bound
barely moving (Gemma-3 4B $0.339\to0.308$, LLaMA-3.2 3B $0.269\to0.251$). This
conclusion does not depend on the size or sign of any single shift: even
granting every improvement at face value, the nearest uncertified model
(Gemma-3 1B, $0.168$) remains well above the $0.10$ target. Because a
perfect risk score would always admit a certifiable threshold ($\alpha_{\min}$
would approach $1/(n+1)$), these failures lie with the scorer or with label
noise, not with an irreducible property of the models: on conversational
queries, Qwen2.5-7B cannot reliably separate regressive from non-regressive
drafts even with the reference answer in hand.

On MedQuAD the two 1B models move to the certified side of the bound (Gemma-3
1B $0.124\to0.079$, LLaMA-3.2 1B $0.120\to0.090$), nominally taking
certification from two of six to four of six. We do not interpret this as a
rescue. Both shifts ($0.045$ and $0.030$) fall below the $0.05$ noise floor
(below), and LLaMA-3.2 1B's is no larger than improvements elsewhere that
change no verdict---Gemma-3 4B improves by $0.031$ on HealthSearchQA yet
stays far from the bound. The verdict flips here only
because these two models already sit near the $0.10$ line, not because
reference access produces a distinctly larger effect, so we report the MedQuAD
crossings as suggestive and rest no conclusion on them.

\begin{table}[t]
\centering
\caption{Oracle diagnostic ($N=1{,}000$, 200 calibration items, both
datasets): $\alpha_{\min}$ with the deployable risk scorer vs.\ one that
observes the reference answer during calibration. Bold: certifies at
$\alpha=0.10$. Not a deployable configuration; reported to localize the
failure. The two additional MedQuAD certifications under the oracle (Gemma-3
1B, LLaMA-3.2 1B) involve $\alpha_{\min}$ shifts within run-to-run noise and
are not interpreted as rescues; see text.}
\label{tab:oracle}
\small
\begin{tabular}{llcccc}
\toprule
& & \multicolumn{2}{c}{MedQuAD} & \multicolumn{2}{c}{HealthSearchQA} \\
\cmidrule(lr){3-4}\cmidrule(lr){5-6}
Family & Size & Deploy. & Oracle & Deploy. & Oracle \\
\midrule
Gemma & 1B   & 0.124 & \textbf{0.079} & 0.187 & 0.168 \\
      & 4B   & 0.184 & 0.201 & 0.339 & 0.308 \\
LLaMA & 1B   & 0.120 & \textbf{0.090} & 0.187 & 0.177 \\
      & 3B   & 0.150 & 0.145 & 0.269 & 0.251 \\
Phi   & 1.3B & \textbf{0.023} & \textbf{0.026} & \textbf{0.081} & \textbf{0.066} \\
      & 2.7B & \textbf{0.048} & \textbf{0.042} & \textbf{0.035} & \textbf{0.028} \\
\midrule
\multicolumn{2}{l}{\emph{certify at $\alpha=0.10$}} & 2/6 & \textbf{4/6} & 2/6 & 2/6 \\
\bottomrule
\end{tabular}
\end{table}

Run-to-run variability is substantial: two independent reruns of this
diagnostic under an identical configuration (same script, seed, and items)
did not reproduce $\alpha_{\min}$ exactly. Both reruns
regenerated every stage---rebuttals, claims, risk scores, and judge
labels---on different GPUs, where greedy decoding is not bit-identical; a
few drafts flip, and because the scorer emits only 12--16 distinct
values per run, $\alpha_{\min}$ moves in discrete jumps. Seven of twelve
model--dataset pairs were identical across reruns and three differed by at
most $0.007$, but the remaining two differed by up to $0.044$, so we treat changes below roughly $0.05$ as
noise. Every per-cell change in Table~\ref{tab:oracle} falls below this
floor---including the two MedQuAD crossings and the two values that rise
slightly (Gemma-3 4B and Phi-1.5 on MedQuAD)---which is why we do not lean on
the MedQuAD crossings, whose verdict flips rest entirely on a sub-noise
movement across the $0.10$ line. The HealthSearchQA reading is immune to this:
it turns on the failure to certify despite the reference answer, not on the
magnitude of any shift.

Two further interventions on the scorer, run as a separate
side experiment, support the same reading. Revising the scoring
prompt left the scorer near or below chance at ranking regressive instances
(AUC $0.302\to0.470$ for LLaMA-3.2 1B and $0.508\to0.575$ for Gemma-3 1B,
measured over 240 calibration instances from a 60-item side run) and did
not change any calibration outcome. An AUC well below $0.5$, as for the
original LLaMA-3.2 1B prompt, indicates a systematically inverted ranking
rather than an uninformative one. Replacing the scorer with a model from a
different family, gpt-oss-20b, gave AUC $0.537$ and $0.418$---
indistinguishable from chance---and calibration still failed, indicating
the failure is not an artifact of the scorer assessing rebuttals generated
by the same model that scores them. Only supplying the reference answer moves any model across the bound, and
those movements fall within run-to-run noise (above). We also observe that the scorer emits only 9--11
distinct score values across those 240 instances, so the threshold scan has
roughly ten candidate operating points; even a well-ranked scorer would
have coarse control over where the accepted region sits.



\section{Discussion}

CCR lowers regressive sycophancy end-to-end for most models, but not
through its conformal components: calibration fails for four of six models,
the resulting universal rewrite is harmful, and the reduction is consistent
with claim filtering, which, where the reduction is largest, acts as an
unplanned abstention mechanism rather than a correction.

Neither susceptibility nor the response to intervention is a simple function of scale. Within each family the larger model is more susceptible at baseline, yet LLaMA-3.2 1B, one of the two smallest models evaluated, is the only one made worse by CCR on both datasets, while Gemma-3 4B shows the largest end-to-end reduction, though mostly through withheld rather than corrected answers. Any deployment assumption that a
safeguard validated on one model transfers to another in the same family, or
that larger models are more robust to such interventions, is unsupported
here.

The oracle diagnostic (\S\ref{sec:oracle}) places the bottleneck in the risk scorer, or in label noise, rather than in the models: even with the reference answer, Qwen2.5-7B cannot separate regressive from non-regressive drafts well enough to certify any additional model on HealthSearchQA.

Finally, CCR never intervenes selectively in any configuration we evaluated. Certified models receive no rewrites and uncertified models receive a rewrite on every response, so these results do not test the premise that rewriting only high-risk responses improves on intervening uniformly.

\section{Limitations}

Our calibration set of 200 examples is modest. At this size the finite-sample
correction is negligible (Section~\ref{sec:config}), but a larger set would
reduce sampling variability in the estimated risk and yield a more stable
threshold. The oracle diagnostic is calibration-only: it reuses the same 200 items and runs no test phase, so it localizes the failure but does not estimate deployed performance. Qwen2.5-7B serves as judge, rebuttal
generator, claim scorer, and risk scorer; a single model in multiple roles
may introduce correlated errors, and the oracle diagnostic suggests
risk-scoring capacity specifically is the limiting factor. Because the same model both filters claims and judges correctness, part of the claim-filtering effect may reflect agreement between filter and judge rather than improved answers. The HealthSearchQA reference answers are themselves model-generated (MedLFQA~\citep{jeong2024olaph}), so label noise on that dataset may be higher than on MedQuAD, whose references are institutionally curated. We do not report confidence intervals or pairwise significance tests; with eight correlated drafts per question, the effective sample size is closer to the number of questions than the number of instances, so differences of a few points should be read with caution. Our definition of regressive sycophancy also excludes transitions from correct to erroneous, so an intervention that degrades answers into unjudgeable output would lower measured regressive sycophancy without improving them. This occurs: Table~\ref{tab:e2e} reports the change in erroneous answers, and where regressive sycophancy falls most, much of the reduction is answers withheld rather than corrected (\S\ref{sec:e2e}). The baseline and CCR runs formatted the in-context conversation history differently for the second to fourth rebuttals, and each run generated its own initial answers; restricting the comparison to the five draft slots with identical prompts changes no sign in Table~\ref{tab:e2e}.
% TODO: add question-level (clustered) CIs and the full correct/incorrect/erroneous transition matrix for baseline, filtered, and rewritten outputs.
 A few MedQuAD questions share identical text; in particular, five test
questions (four distinct texts) share their text with calibration questions. Because the baseline
and intervention arms are paired by question hash, the drafts of these five
questions are dropped from the baseline arm, which therefore contains
$6{,}360$ rather than $6{,}400$ instances. We did not deduplicate the split;
this affects under 1\% of test questions and has a negligible effect on the
reported rates, but it means a small number of test questions also appear
in the calibration set. Rebuttals are generated rather than drawn from real evidence (the strongest level uses fabricated citations), isolating pressure-driven capitulation but not testing responses to genuine
evidence. Finally, evaluation is limited to open-source models of at most
4B parameters; larger and RLHF-aligned models, where sycophancy may be more
deeply embedded, are not evaluated.

\section{Conclusion}

We evaluated conformal risk control as a deployable post-hoc safeguard
against sycophancy in medical question answering, under the constraint that
no component acting on a response at test time may access a reference
answer. Using a finite-sample bound, only two of six models admit a
certifiable threshold at a 10\% target---the two base Phi models, whose
low initial accuracy makes the bound easy to meet---and the same two certify
on both datasets. Where calibration fails, the resulting universal rewrite increases
regressive sycophancy in every case we measured. The reduction observed
end-to-end is consistent with claim-level filtering operating at an
uncalibrated fallback threshold rather than with the conformal components,
and where it is largest it consists mainly of answers withheld rather than
corrected. An oracle diagnostic indicates that the binding constraint is the risk scorer: even with the reference answer supplied during calibration, no model on HealthSearchQA becomes certifiable, and the only verdict changes, on MedQuAD, are within run-to-run noise. Conformal
calibration is valuable here less as a mitigation than as an audit
instrument: it detects, using a finite-sample criterion, when LLM-based risk
scoring cannot support selective intervention, and $\alpha_{\min}$
quantifies the gap any such method must close.

% \begin{acks}
% Add acknowledgments for camera-ready.
% \end{acks}



\section*{Ethics and Privacy Statement}

This work uses only public datasets (MedQuAD, and the HealthSearchQA questions with the publicly released MedLFQA reference answers) that contain no patient or protected health information, and involves no human subjects. Adversarial rebuttals, including fabricated citations, were generated solely to probe model behavior and are not released as standalone content. CCR is a research diagnostic, not a clinical tool: our results show that it does not reliably reduce sycophancy through its calibrated components, and a passing conformal certificate should not be read as evidence that a model is safe for medical use.

\section*{GenAI Usage Disclosure}

Generative AI tools, including Anthropic's Claude and Claude Code, were
used throughout this work. Specifically, they were used to (i) search for
and summarize related literature and format bibliography entries; (ii)
write and debug the experiment and analysis code; (iii) generate figures
and tables; (iv) produce initial drafts of manuscript text; and (v) review
drafts, identify internal inconsistencies, suggest and apply wording edits,
and inspect result files to compute descriptive checks reported in the paper
(e.g., rate denominators and initial-accuracy statistics). All AI-assisted
text, code, figures, and analyses were reviewed and verified by the
authors, who take full responsibility for the content of this paper.
Separately, large language models are the object of study: the evaluated
models and the auxiliary Qwen2.5-7B-Instruct and gpt-oss-20b components
described in Sections~\ref{sec:method} and~\ref{sec:oracle} are part of the
experimental method, not writing aids.

\bibliographystyle{ACM-Reference-Format}
\bibliography{sample-base}

\end{document}
\endinput
%%
%% End of file `sigconf.tex'.
